\chapter{Covering arguments for singular strata}\label{sec:strata-covering}
Let $V$ be a stationary integral $m$-varifold in $U\subset\mathbb R^N$, $1\leq m<N$. Write $\Theta(x)=\Theta^m(\|V\|,x)$ and use the strata $\mathcal S^j(V)$ of Definition \ref{def:singular-strata}. The argument below expands \cite{Minter}, Theorem 1.7 and its proof, pp.~17--19 and 22--24. It uses every possible tangent cone through the definition of the strata, but never assumes tangent-cone uniqueness.

\begin{definition}[One-sided approximation by planes]\label{def:one-sided-plane-approximation}
A set $A\subset\mathbb R^N$ has the $(j,\delta,r_0)$ approximation property if for every $a\in A$ and $0<r\leq r_0$ there is a $j$-dimensional linear subspace $L=L(a,r)$ such that
\[
 A\cap B_r(a)\subset\{x:\operatorname{dist}(x,a+L)<\delta r\}.
\]
Here $0\leq j<N$ and $\delta,r_0>0$. The plane may depend on both center and scale; no converse approximation of the plane by $A$ is required. Compare \cite{Minter}, p.~19.
\end{definition}

\begin{lemma}[Refining a cover near a plane]\label{lem:flat-cover-refinement}
There is $C_N\geq1$ such that, if $0<\delta<1/16$ and $A$ has the $(j,\delta,r_0)$ approximation property, each $A\cap B_r(a)$, $a\in A$, $r\leq r_0$, can be covered by at most $C_N\delta^{-j}$ balls of radius $8\delta r$ whose centers belong to $A\cap B_r(a)$.
\end{lemma}
\begin{proof}
Project orthogonally onto the approximating plane $a+L$. The projected points lie in its $j$-ball of radius $r$. A cubical grid covers this ball by at most $C_N\delta^{-j}$ cells of diameter at most $\delta r$. For each cell whose inverse image meets $A\cap B_r(a)$, select one point of that intersection. Two points over the same cell have projected distance at most $\delta r$ and normal components of length less than $\delta r$ each; their distance is less than $3\delta r$. The balls of radius $8\delta r$ at the selected points therefore cover the set. If $j=0$, the projection is a single point and one such ball suffices. Centers remain in $A$, which allows iteration of the approximation hypothesis.
\end{proof}

\begin{proposition}[Small one-sided approximation bounds Hausdorff dimension]\label{prop:flat-approximation-dimension}
For every $\beta>0$ there is $\delta_*(N,j,\beta)>0$ such that any set with the $(j,\delta,r_0)$ approximation property, $0<\delta<\delta_*$, has
\[
 \mathcal H^{j+\beta}(A)=0.
\]
No measurability or closedness assumption on $A$ is needed for this outer-measure conclusion.
\end{proposition}
\begin{proof}
Set $s=j+\beta$ and $\vartheta=8\delta$. Choose $\delta_*<1/16$ so small that
$C_N\delta^{-j}\vartheta^s<1/2$ whenever $0<\delta<\delta_*$. This is possible because the expression is $C_N8^s\delta^\beta$. First suppose $A$ is bounded. Total boundedness supplies a finite cover by balls of a fixed radius $r<r_0$ with centers in $A$. Let the sum of their diameter-to-the-$s$ powers be $T$.

Apply Lemma \ref{lem:flat-cover-refinement} to every ball and iterate. Each new center still lies in $A$ and each radius is smaller than $r_0$. At the $k$th stage all radii are $\vartheta^kr$ and the sum of diameter-to-the-$s$ powers is at most $2^{-k}T$. These are arbitrarily fine covers with cost tending to zero, so $\mathcal H^s(A)=0$, regardless of its normalization constant. For unbounded $A$, apply the argument to $A\cap B_\ell(0)$, which inherits the same property, and take a countable union.
\end{proof}

\begin{lemma}[Nearly maximal density concentrates near a spine]\label{lem:stratum-density-concentration}
For $a\in\mathcal S^j(V)$ and $\delta>0$, there is $\varepsilon>0$ with $\overline{B_{2\varepsilon}(a)}\subset U$ such that for each $0<r\leq\varepsilon$ some $j$-plane $L=L(a,r)$ satisfies
\[
 \{y\in B_r(a):\Theta(y)\geq\Theta(a)-\varepsilon\}
 \subset\{y:\operatorname{dist}(y,a+L)<\delta r\}.
\]
\end{lemma}
\begin{proof}
If the assertion fails, choose $\varepsilon_i\downarrow0$ and $r_i\leq\varepsilon_i$ for which no $j$-plane works. Theorem \ref{thm:tangent-cones} gives a subsequence of $V_{a,r_i}$ converging locally to a tangent cone $C$. Its spine has dimension at most $j$ by the definition of $\mathcal S^j$. Choose a $j$-plane $L$ containing this spine. Failure for this particular plane gives points $z_i\in B_1(0)$ with
\[
 \Theta^m(\|V_{a,r_i}\|,z_i)\geq\Theta(a)-\varepsilon_i,
 \qquad \operatorname{dist}(z_i,L)\geq\delta.
\]
Pass to $z_i\to z\in\overline{B_1(0)}$. Density upper semicontinuity, Proposition \ref{prop:density-upper-semicontinuity}, gives $\Theta^m(\|C\|,z)\geq\Theta(a)=\Theta^m(\|C\|,0)$. Proposition \ref{prop:spine-maximal-density} forces $z\in S(C)\subset L$, contradicting the distance bound. Convergence is on all compact subsets of $\mathbb R^N$, so a limit with $|z|=1$ causes no boundary problem.
\end{proof}

\begin{proposition}[Decomposing a stratum by scale and density]\label{prop:stratum-flat-pieces}
For every fixed $\delta>0$, $\mathcal S^j(V)$ is a countable union of sets, each having the $(j,\delta,1/i)$ approximation property for some positive integer $i$.
\end{proposition}
\begin{proof}
Let $E_i$ consist of points $a\in\mathcal S^j(V)$ for which the conclusion of Lemma \ref{lem:stratum-density-concentration} holds with tolerance and scale $\varepsilon=1/i$, and $\overline{B_{2/i}(a)}\subset U$. Every stratum point belongs to some $E_i$: decreasing the tolerance and allowed scale preserves the conclusion. Subdivide into
\[
 E_{i,q}=\{a\in E_i:(q-1)/i<\Theta(a)\leq q/i\},\qquad q=1,2,\ldots.
\]
Densities are finite and at least one at support points, so these sets cover the stratum. If $a,y\in E_{i,q}$, then $\Theta(y)>\Theta(a)-1/i$. Thus the concentration conclusion at $a$, for any $r\leq1/i$, applies to every $y\in E_{i,q}\cap B_r(a)$. This is precisely the desired approximation property.
\end{proof}

\begin{corollary}[The zero stratum is countable]\label{cor:zero-stratum-countable}
The set $\mathcal S^0(V)$ is at most countable.
\end{corollary}
\begin{proof}
Use Proposition \ref{prop:stratum-flat-pieces} with a fixed $0<\delta<1$. A piece with the $(0,\delta,r_0)$ property has no distinct points $a,b$ with $d=|a-b|<r_0$: choose $d<r\leq r_0$ also satisfying $\delta r<d$. The property at $a$ would give $b\in B_{\delta r}(a)$, a contradiction. Each piece is therefore $r_0$-separated and has finitely many points in each bounded region. It is countable, as is the union of pieces. No local finiteness of the entire zero stratum is asserted.
\end{proof}

\begin{definition}[Flat singular points]\label{def:flat-singular-point}
A singular point of $V$ is called flat if at least one tangent cone there is a positive-integer multiple of an $m$-plane. The point is still required to be singular according to Definition \ref{def:varifold-regular-singular}; a regular multiple plane is not a flat singular point. This terminology follows \cite{Minter}, p.~20.
\end{definition}

\begin{corollary}[Flat singular points form the top stratum difference]\label{cor:flat-singular-top-stratum}
The flat singular points are exactly $\mathcal S^m(V)\setminus\mathcal S^{m-1}(V)$. At each such point the density is an integer $q\geq2$.
\end{corollary}
\begin{proof}
Failure of membership in $\mathcal S^{m-1}$ means some tangent cone has spine dimension $m$, since Proposition \ref{prop:spine-maximal-density} bounds it above by $m$. Corollary \ref{cor:full-spine-plane} identifies that cone as $q|P|$. Conversely a multiple plane has spine dimension $m$. The tangent-cone mass formula, Theorem \ref{thm:tangent-cones}, identifies the density at the original point as $q$. If $q=1$, Corollary \ref{cor:unit-density-regularity} makes the point regular, a contradiction.
\end{proof}

\begin{corollary}[The singular set below density two]\label{cor:subtwo-singular-dimension}
One has
\[
 \dim_{\mathcal H}\{a\in\operatorname{sing}V:\Theta(a)<2\}\leq m-1.
\]
For $m=1$ this set is at most countable.
\end{corollary}
\begin{proof}
Corollary \ref{cor:flat-singular-top-stratum} excludes the top stratum difference under the density bound. The set is therefore contained in $\mathcal S^{m-1}(V)$. Apply Theorem \ref{thm:strata-dimension}; for $m=1$ use the stronger conclusion of Corollary \ref{cor:zero-stratum-countable}. This does not assert regularity merely from density less than two.
\end{proof}
